\section{Evaluations filtering}
\label{app:results-filtering}

For each benchmark--language subset $(B,L)$ and model size $N$, we test whether the observed accuracy is distinguishable from random guessing.
For a task with $K$ answer options, chance accuracy is defined as $1/K$.
For every evaluation run, we compute the one-sided 95\% Wilson lower confidence bound on its accuracy and mark the run as above chance when

\[
\operatorname{LCB}_{95}(\hat{p}) > \frac{1}{K}.
\]

We then aggregate these decisions across the available architectures, language mixtures, and seeds for the same $(B,L,N)$ cell.
A cell is retained if at least 50\% of its evaluation runs are above chance.
When models trained on the benchmark language are available, we use those runs for the filtering decision; otherwise, we use all evaluated runs.
Measurements without a meaningful chance baseline, such as bits per byte and training loss, are not filtered by this procedure.

The Wilson bound accounts for the number of evaluation examples, so the amount by which accuracy must exceed chance varies across tasks.
Smaller evaluation sets require a larger margin above chance, while larger sets can support smaller but more reliable departures from chance.
For example, the median required margins are 9.00 percentage points for nonparallel Global PIQA, 2.44 for Belebele, 0.74 for HellaSwag, and 0.60 for Global-MMLU.
The corresponding fractions of benchmark--language--size cells that pass the filter are 37.5\%, 6.8\%, 92.7\%, and 4.3\%, respectively.
This filtering step is used only to determine whether a benchmark provides measurable signal at a given scale; it does not assess whether the benchmark can reliably rank model variants.
